% Chapter 2 — includable by main.tex
\chapter{Regularni izrazi}
\label{chap:2}

\begin{quote}
\itshape
U prvom poglavlju naučili smo što su jezici i kako ih kombinirati.
Sad postavljamo pitanje: kako ih \emph{opisati}? Regularni izrazi daju
nam kompaktan, elegantan formalizam koji bi trebao izgledati poznato ---
vjerojatno ste već koristili \texttt{grep} ili \texttt{re} modul u Pythonu.
Ovo poglavlje tu intuiciju pretvara u preciznu matematiku.
\end{quote}


% ================================================================
\section{Motivacija: kako opisati beskonačni jezik?}
% ================================================================

Zamislite da vam netko kaže: ,,napiši mi sve valjane e-mail adrese''.
Odmah je jasno da taj skup ne možete nabrojati --- beskonačni je i nema
gornje granice duljine. No ipak intuitivno znate što je valjana adresa:
neka slova, pa \texttt{@}, pa domena. Regularni izrazi su upravo alat koji
takav opis pretvara u preciznu matematičku definiciju.

U Pythonu ste možda pisali nešto poput \verb|re.match(r'[a-z]+@[a-z]+\.[a-z]+', s)|.
Formalni regularni izrazi koje ćemo definirati su strožija i čišća verzija
iste ideje --- bez šarenih proširenja koja praktični jezici dodaju, ali s
jasnom matematičkom semantikom.

Ideja je jednostavna: regularni izrazi grade se od malih komadića
(\emph{inicijalnih} ili \emph{atomarnih} jezika) kombiniranjem s tri
operacije --- unije, konkatenacije i Kleeneve zvijezde. Ništa više.
Pokazat će se da je to i više nego dovoljno za cijelu klasu regularnih
jezika.

% ================================================================
\section{Definicija regularnih izraza}
% ================================================================

\begin{definicija}[RI-regularni jezik]
Neka je $\Sig$ abeceda. \textbf{Skup RI-regularnih jezika} nad $\Sig$
je \emph{najmanji} skup jezika koji:
\begin{itemize}
  \item sadrži $\emptyset$ i $\{\eps\}$ kao elemente,
  \item sadrži $\{\alpha\}$ za svaki znak $\alpha \in \Sig$,
  \item je zatvoren na konkatenaciju, konačnu uniju i Kleeneve operacije.
\end{itemize}
Jezik koji pripada tom skupu zovemo \textbf{RI-regularnim}.
\textbf{Regularni izraz} je izraz koji pokazuje kako je taj jezik izgrađen
od navedenih operacija.
\end{definicija}

Dakle, regularni izraz nije sam po sebi jezik --- on je \emph{recept} za
gradnju jezika. Kao što izraz $2 \cdot 3 + 1$ nije broj već opis kako
doći do broja $7$, regularni izraz $ab^*$ nije skup nego opis kako
doći do skupa $\{a, ab, abb, abbb, \ldots\}$.

\subsection{Inicijalni jezici}

Gradivni blokovi su sljedeća tri tipa:

\begin{center}
\renewcommand{\arraystretch}{1.5}
\begin{tabular}{cll}
\hline
\textbf{Izraz} & \textbf{Jezik} & \textbf{Opis} \\
\hline
$\emptyset$ & $\emptyset$ & prazan jezik, nema nijedne riječi \\
$\eps$      & $\{\eps\}$  & samo prazna riječ \\
$\alpha$    & $\{\alpha\}$ & samo jednoznakovna riječ od $\alpha$ \\
\hline
\end{tabular}
\end{center}

\begin{napomena}
Pažnja na oznake: $\emptyset$ kao regularni izraz opisuje \emph{prazan jezik},
a $\eps$ kao regularni izraz opisuje jezik $\{\eps\}$ koji sadrži \emph{jednu}
riječ. Ovo je ista razlika kao u prethodnom poglavlju --- samo sad gledamo
izraze, ne skupove.
\end{napomena}

\subsection{Operacije}

Iz inicijalnih jezika gradimo veće regularnim operacijama:

\begin{center}
\renewcommand{\arraystretch}{1.5}
\begin{tabular}{cll}
\hline
\textbf{Izraz} & \textbf{Jezik} & \textbf{Čitamo} \\
\hline
$r \cup s$  & $L(r) \cup L(s)$   & unija: ili $r$ ili $s$ \\
$r \cdot s$ & $L(r) \cdot L(s)$  & konkatenacija: prvo $r$, pa $s$ \\
$r^*$       & $L(r)^*$           & zvijezda: nula ili više puta $r$ \\
$r^+$       & $L(r)^+$           & plus: jedanput ili više puta $r$ \\
\hline
\end{tabular}
\end{center}

Točka za konkatenaciju često se izostavlja: $rs$ znači isto što i $r \cdot s$.
Prioritet operacija je kao u aritmetici: ${}^*$ i ${}^+$ vežu najjače, pa
konkatenacija, pa unija. Dakle $ab^*c \cup d$ čitamo kao
$(a \cdot (b^*) \cdot c) \cup d$.

% ================================================================
\section{Primjeri --- od jednostavnog prema složenom}
% ================================================================

Sad kad imamo definiciju, vježbamo je na konkretnim primjerima. Ključna
vještina je dvojaka: čitati regularni izraz i razumjeti koji jezik opisuje,
i zadani jezik opisati regularnim izrazom.

\begin{primjer}[Jednostavni izrazi]
Nad $\Sig = \{a, b\}$:
\begin{itemize}
  \item $a \cup b$ opisuje jezik $\{a, b\}$ --- jedna od ta dva slova.
  \item $ab$ opisuje $\{ab\}$ --- točno ta jedna dvoznakovna riječ.
  \item $a^*$ opisuje $\{\eps, a, aa, aaa, \ldots\}$ --- nula ili više $a$-ova.
  \item $a^+$ opisuje $\{a, aa, aaa, \ldots\}$ --- jedan ili više $a$-ova.
  \item $a^* b$ opisuje $\{b, ab, aab, aaab, \ldots\}$ --- nula ili više $a$-ova, pa jedan $b$.
\end{itemize}
\end{primjer}

\begin{primjer}[Sve riječi nad $\{a, b\}$]
Izraz $(a \cup b)^*$ opisuje \emph{sve} moguće riječi nad $\{a, b\}$,
dakle $\Sig^*$. Kleeneva zvijezda od $\{a\} \cup \{b\} = \{a,b\}$
daje sve načine nizanja $a$-ova i $b$-ova, bez ikakva ograničenja.

To je korisna crtica: kad god trebate ,,bilo što od tih znakova'',
pišete $(a \cup b)^*$, ili općenito $\Sig^*$.
\end{primjer}

\begin{primjer}[Binarni brojevi]
Razmotrimo izraz $1(0 \cup 1)^* \cup 0$ nad $\Sig = \{0, 1\}$.
\begin{itemize}
  \item $1(0 \cup 1)^*$: počinje s $1$, a zatim slijedi bilo što ---
        to su svi binarni brojevi $\geq 1$.
  \item $0$: baš sama nula.
  \item Unija: skupa dobivamo \emph{sve valjane binarne zapise prirodnih
        brojeva} bez vodećih nula (jedino $0$ počinje s nulom).
\end{itemize}
Ovo je tipičan uzorak: posebnim slučajem ($0$) riješimo rub,
a glavnom granom ($1(0\cup1)^*$) pokrijemo opći slučaj.
\end{primjer}

\begin{primjer}[Riječi koje počinju i završavaju istim slovom]
Napišimo regularni izraz za sve riječi nad $\{a, b, c\}$ koje
počinju i završavaju istim slovom. Koristimo kraticu $\Sig := a \cup b \cup c$.

Fiksiramo koje je to slovo i ujedinimo tri slučaja:
\begin{itemize}
  \item Počinje i završava s $a$: treba $a \Sig^* a$, ali i samu riječ $a$.
        Dakle: $a(\Sig^* a \cup \eps)$.
  \item Analogno za $b$ i $c$.
\end{itemize}
Ukupno:
\[
  a(\Sig^* a \cup \eps) \;\cup\; b(\Sig^* b \cup \eps) \;\cup\; c(\Sig^* c \cup \eps).
\]
\end{primjer}

\begin{primjer}[Neparan broj $a$-ova]
Jezik svih riječi nad $\{a,b\}$ s neparnim brojem $a$-ova.

Intuicija: između $a$-ova može biti koliko god $b$-ova. Pokušajmo
$b^* (ab^*)^+$ --- ali to dopušta \emph{bilo koji} broj $a$-ova, parni ili neparni.

Rješenje: svaki par $a$-ova se ,,poništi'', a jedan ostaje:
\[
  b^* \bigl(ab^* ab^*\bigr)^* ab^*.
\]
Čitanje: nula ili više parova $a \cdots a$ (okruženih $b$-ovima),
pa točno jedan $a$, pa $b$-ovi. Ukupan broj $a$-ova je $2k + 1$ za neko
$k \geq 0$ --- uvijek neparan.
\end{primjer}

\begin{kljucno}
Izgradnja regularnog izraza za zadani jezik obično ide ovako:
\begin{enumerate}
  \item Identificirajte strukturu riječi --- što dolazi na kojoj poziciji.
  \item Podijelite na slučajeve ako treba (npr.\ različiti početni znakovi).
  \item Koristite $(a \cup b \cup \ldots)^*$ za ,,bilo što od tih znakova''.
  \item Provjerite rubne slučajeve: prazna riječ, jednoznakovne riječi.
\end{enumerate}
\end{kljucno}

% ================================================================
\section{Algebarske jednakosti}
% ================================================================

Dva regularna izraza $r$ i $s$ su \textbf{ekvivalentni} ako opisuju isti
jezik: $L(r) = L(s)$. Pišemo $r \equiv s$.

Postoji niz korisnih algebarskih identiteta:

\begin{center}
\renewcommand{\arraystretch}{1.6}
\begin{tabular}{ll}
\hline
\textbf{Identitet} & \textbf{Što kaže} \\
\hline
$r \cup \emptyset \equiv r$ & unija s praznim jezikom ne mijenja ništa \\
$r \cdot \eps \equiv r$     & konkatenacija s $\eps$ ne mijenja ništa \\
$r \cdot \emptyset \equiv \emptyset$ & konkatenacija s praznim daje prazno \\
$r \cup r \equiv r$         & unija sa samim sobom je idempotentna \\
$r \cup s \equiv s \cup r$  & unija je komutativna \\
$(rs)t \equiv r(st)$        & konkatenacija je asocijativna \\
$r(s \cup t) \equiv rs \cup rt$ & distributivnost \\
$(r^*)^* \equiv r^*$        & zvijezda od zvijezde je zvijezda \\
$\eps^* \equiv \eps$        & zvijezda od $\eps$ je $\eps$ \\
$\emptyset^* \equiv \eps$   & zvijezda od praznog skupa je $\{\eps\}$ \\
\hline
\end{tabular}
\end{center}

\begin{primjer}
Provjerimo $\emptyset^* \equiv \eps$, tj.\ $L(\emptyset^*) = \{\eps\}$.
Po definiciji, $\emptyset^* = \bigcup_{k \geq 0} \emptyset^k$.
Vrijedi $\emptyset^0 = \{\eps\}$ i $\emptyset^k = \emptyset$ za $k \geq 1$.
Dakle $\emptyset^* = \{\eps\}$. Intuitivno: jedini način da se ,,niže nula
ili više riječi iz praznog skupa'' je uzeti nula riječi, što daje $\eps$.
\end{primjer}

\begin{zadatak}
\begin{enumerate}[label=(\alph*), itemsep=4pt]
  \item Pokažite da je $r^+ \equiv rr^* \equiv r^*r$.
  \item Vrijedi li $r^* \equiv r^+ \cup \eps$ za svaki $r$?
        Što ako je $\eps \in L(r)$?
  \item Nad $\Sig = \{0,1\}$, napišite regularni izraz za:
    \begin{itemize}
      \item sve riječi koje završavaju na $01$,
      \item sve riječi koje sadrže podniz $11$,
      \item sve riječi parne duljine.
    \end{itemize}
  \item Nad $\Sig = \{a, b\}$, napišite regularni izraz za sve riječi
        u kojima se $a$ nikad ne pojavljuje dva puta zaredom.
\end{enumerate}
\end{zadatak}

% ================================================================
\section{Što regularni izrazi ne mogu?}
% ================================================================

Regularni izrazi su moćni, ali nisu svemoćni.

Jezik $L_= = \{a^n b^n \mid n \geq 1\}$ sadrži riječi $ab$, $aabb$,
$aaabbb$, \ldots\ Svaka riječ ima isti broj $a$-ova i $b$-ova.
Pokušajte napisati regularni izraz za taj jezik --- neće ići.

Zašto? Regularni izraz ne može ,,prebrojati'' i usporediti duljine dva
dijela riječi. Sve što može je pratiti konačno mnogo mogućih ,,stanja''
--- a mogućih vrijednosti $n$ je beskonačno mnogo. Formalni dokaz dolazi
s \textbf{lemom o pumpanju} (poglavlje~6), ali intuicija je jasna već sad.

\begin{kljucno}
Regularni izrazi opisuju jezike u kojima se struktura može pratiti bez
pamćenja neograničenih brojača. Čim trebate ,,zapamtiti'' koliko puta
ste nešto vidjeli i usporediti s nekim drugim brojem --- regularni izraz
nije dovoljan alat.
\end{kljucno}

% ================================================================
\section*{Sažetak}
% ================================================================

Regularni izrazi daju nam kompaktan način opisivanja regularnih jezika.
Grade se od tri atomarna slučaja ($\emptyset$, $\eps$, jednoznakovni jezici)
i tri operacije (unija, konkatenacija, Kleeneva zvijezda/plus).

\medskip
\noindent
\begin{tabular}{ll}
  $\emptyset$  & prazan jezik \\
  $\eps$       & jezik $\{\eps\}$ \\
  $\alpha$     & jezik $\{\alpha\}$ \\
  $r \cup s$   & unija \\
  $rs$         & konkatenacija \\
  $r^*$        & Kleeneva zvijezda (nula ili više puta) \\
  $r^+$        & Kleenejev plus (jedanput ili više) \\
\end{tabular}

\medskip
\noindent U sljedećem poglavlju uvodimo \textbf{konačne automate} ---
strojeve koji prepoznaju regularne jezike. Vidjet ćemo da su automati i
regularni izrazi \emph{ekvivalentni}: svaki regularni izraz daje automat,
i svaki automat daje regularni izraz.
